\documentclass{situationist-preprint}

\title{Projectibility and Naturalness}
\msid{EPI-04-a}
\msfamily{EPI-04}
\msdiscipline{Epistemology}
\mstype{Principal manuscript}
\mskeywords{new riddle of induction; grue; projectibility; natural properties; sparse universals; entrenchment; lawlikeness}

\begin{document}
\maketitle

\begin{abstract}
All emeralds examined so far have been green. They have also been grue, where an object is grue if it is green and examined before a given future time, or blue and not so examined. The same evidence seems to confirm that all emeralds are green and that all emeralds are grue, which make incompatible predictions. Goodman's new riddle asks why \enquote{green} is projectible and \enquote{grue} is not. We argue that \enquote{green} denotes a natural property and \enquote{grue} does not, and that projectibility tracks naturalness. Relative naturalness is measured by the complexity of the shortest construction of a property from the sparse universals. Inductive inference projects resemblance, objective resemblance consists in the sharing of natural properties, and the regularities that hold by law are stated in natural terms; so, given fixed evidence and a fixed kind, the more natural of two predicates is the more projectible. Grue, unlike green, distinguishes intrinsic duplicates, and so can be constructed only with a particular time as a parameter, by a construction that contains a construction of green. We answer the symmetry objection, the charge of circularity and the claim that entrenchment suffices, and examine cases in which time-dependent properties are projectible.
\end{abstract}

\section{The problem}

\citet{Goodman1946,Goodman1955} posed what he called the new riddle of induction. Let $t$ be a time in the future. Say that an object is \emph{grue} if it is green and examined before $t$, or blue and not examined before $t$. Every emerald examined so far has been green, and, since it was examined before $t$, grue. If an instance of a generalization confirms the generalization, as the instance theory of confirmation of \citet{Hempel1945} holds, the evidence confirms both that all emeralds are green and that all emeralds are grue. The two hypotheses agree about examined emeralds and disagree about every emerald first examined after $t$: the first predicts that it will be green, the second that it will be blue. We take the first to be confirmed and the second not, and the riddle is to say why. Following Goodman, call a predicate \emph{projectible} if its observed instances confirm that unobserved members of the relevant kind share it.

The riddle is not a puzzle about one predicate. For any body of evidence there are indefinitely many grue-like predicates that fit it and diverge on any unobserved case, so that a theory of confirmation that does not discriminate among predicates confirms every prediction equally. EPI-03 argues that induction is justified, inductively, and observes that the inductive rule \enquote{predict that the unobserved resembles the observed} is well defined only once the respects of resemblance are fixed; it relies on the present manuscript to fix them. The answer to the new riddle is therefore part of the answer to the old one.

The main responses are four. \citet{Carnap1947} required that the predicates of an inductive logic be purely qualitative, containing no reference to particular times, places or objects. Goodman replied that qualitativeness is relative to a language: define \emph{bleen} as blue and examined before $t$, or green and not; then \enquote{green} is definable as grue and examined before $t$ or bleen and not, and in the grue--bleen language it is \enquote{green} that refers to a time. Goodman's own answer was \emph{entrenchment}: a predicate is projectible when it, or a predicate coextensive with it, has been projected more often in the past \citep[ch.~IV]{Goodman1955}. \citet{Quine1969} held that projectible predicates are those that correspond to natural kinds, and that our innate similarity standards, refined by science, track them. Bayesian treatments hold that the riddle is answered by prior probabilities, which favour the green hypothesis over the grue one \citep{Sober1994}, and some have concluded that the riddle shows only that no inductive inference proceeds without substantive background assumptions \citep{Okasha2007,Norton2006,GodfreySmith2003}.

We defend a version of Quine's answer, with naturalness in place of natural kinds and with a measure of naturalness that does not depend on any language. Our thesis is that \enquote{green} is projectible and \enquote{grue} is not because \enquote{green} denotes a natural property and \enquote{grue} does not, and that, more generally, projectibility tracks naturalness.

\section{Framework}

We take over the theory of properties defended in MET-02. There are few genuine universals, and every local qualitative property is constructed from them by negation, conjunction, disjunction and projection along genuine relations (MET-02-a, Definition 3.4). Properties that are not qualitative, such as being two miles from this sphere, are constructed from universals with particular objects as parameters.

\begin{definition}[Relative naturalness]\label{d:nat}
The \emph{naturalness} of a property varies inversely with the complexity of its shortest construction from the sparse universals, counting particular objects used as parameters as part of the construction. The genuine universals are the perfectly natural properties.
\end{definition}

Definition~\ref{d:nat} is the measure proposed by \citet[p.~61]{Lewis1986}, on which the less natural properties are those definable from the perfectly natural ones only by long chains of definition, and it is the measure adopted in MET-02-a. Complexity is measured from the universals, not from the predicates of any language. Two languages that differ in their primitives may differ in how long a predicate's definition is, but they do not differ in how long a property's construction from universals is.

\begin{definition}[Projectibility]\label{d:proj}
A predicate $F$ is \emph{projectible} over a kind $K$ to the degree that the evidence that observed $K$s are $F$ confirms that unobserved $K$s are $F$.
\end{definition}

Projectibility so defined is comparative: a predicate is more or less projectible over a kind, given evidence of a given size. The thesis that projectibility tracks naturalness is the following.

\begin{principle}[Tracking]\label{p:track}
Given a kind $K$ and evidence that the same observed $K$s are $F$ and are $G$, $F$ is more projectible over $K$ than $G$ if and only if $F$ is more natural than $G$, and equally projectible if they are equally natural.
\end{principle}

\ref{p:track} holds fixed the kind and the evidence. It does not say that every natural predicate is projectible over every kind: being green is not projectible over pebbles, because pebbles do not share a colour in virtue of their kind. It says that, among the hypotheses that fit a body of evidence about a kind, the more natural are the better confirmed.

We use three premises.

\begin{principle}[Resemblance]\label{p:resemblance}
Inductive inference projects resemblance: from the observed $K$s' being alike in a respect, it infers that the unobserved $K$s are alike with them in that respect.
\end{principle}

\begin{principle}[Natural resemblance]\label{p:natres}
Objective resemblance consists in the sharing of natural properties, and two objects resemble one another in a respect to the degree that the property they share in that respect is natural.
\end{principle}

\begin{principle}[Lawlikeness]\label{p:law}
A generalization over a kind is lawlike to the degree that it is a theorem of the best system, or follows from it together with conditions satisfied by the kind; and the axioms of the best system are stated in terms of natural properties.
\end{principle}

\ref{p:resemblance} is the content of the inductive rule of EPI-03. \ref{p:natres} is the connection between resemblance and naturalness defended by \citet{Lewis1983} and proved, for the genuine properties, in MET-02-a: two objects resemble one another in virtue of sharing a genuine property, and sharing a property that is a long construction from the genuine ones is a correspondingly weaker resemblance. \ref{p:law} is the best-system account of laws defended in MET-12, with the requirement, due to \citet{Lewis1983}, that the system be stated in natural terms; without it, a single axiom \enquote{everything is $F$}, for a suitably gerrymandered $F$, would be maximally simple and maximally strong.

\section{The argument}

\subsection{Two routes to tracking}

The first route goes through resemblance. By \ref{p:resemblance}, to project $F$ from observed to unobserved $K$s is to infer that the unobserved $K$s resemble the observed ones in respect $F$. By \ref{p:natres}, sharing a property is a resemblance to the degree that the property is natural. If $F$ is a long construction from the universals, two objects that share it may differ in every natural respect and resemble one another hardly at all. The evidence that observed $K$s share $F$ is then evidence of very little resemblance among them, and gives correspondingly little reason to expect the unobserved to share it. If $F$ is perfectly natural, the observed $K$s resemble one another in a genuine respect, and the resemblance is what the inductive rule projects. So, given the same evidence, the more natural predicate is projected more strongly by the rule that EPI-03 shows to be justified in a uniform world.

The second route goes through laws. Observed instances of a generalization confirm it, in the sense relevant to prediction, when they are evidence that it holds by law or follows from laws, and not merely that it has held so far. This is the distinction between lawlike and accidental generalizations that Goodman took the riddle to illuminate \citep[ch.~III]{Goodman1955}: that all the coins in my pocket are silver is not confirmed, for the next coin to be put there, by the silver coins already there. By \ref{p:law}, a generalization is lawlike to the degree that it is stated in, or derivable from, natural terms. A generalization in terms of a predicate with a long construction can follow from the best system only through the conditions that the construction encodes, and each of these is a further contingency that the observed instances do not confirm. So, again, the more natural of two predicates that fit the same evidence yields the better-confirmed projection.

The two routes converge because they are two aspects of a single fact. The best system is stated in natural terms because natural properties are the respects in which things objectively resemble, and the inductive rule tracks resemblance because, in a world governed by a system stated in those terms, resemblance in natural respects is what persists from observed to unobserved cases. \ref{p:track} follows from either route; the second explains why the first works.

\subsection{Green and grue}

It remains to show that green is more natural than grue. Green, or the surface reflectance profile that realizes it, is not a genuine universal; it is a constructed property, and its shortest construction from the universals of physics is long. Grue is constructed too. The question is which construction is shorter.

\begin{proposition}\label{p:dup}
Grue distinguishes intrinsic duplicates; green does not.
\end{proposition}

\begin{proof}
Let $a$ and $b$ be emeralds that are intrinsic duplicates, alike in every intrinsic genuine property, and suppose $a$ is examined before $t$ and $b$ is first examined after $t$. Both are green, since green is a local qualitative property and they agree on every genuine property on which local qualitative properties supervene (MET-02-a, Principle P4 and Theorem 5.6). But $a$ is grue, being green and examined before $t$, and $b$ is not, being green and not examined before $t$.
\end{proof}

By Proposition~\ref{p:dup}, grue is not a local qualitative property: it does not supervene on the intrinsic genuine properties of its bearers. Its construction must therefore go beyond them, projecting along genuine relations to other things: to the events of examination, which relate the emerald to observers, and to the particular time $t$, which can enter a construction only as a parameter. Grue is, in the terminology of MET-02-a, an impure property, constructed from universals with a particular as a parameter. And since grue coincides with green on everything examined before $t$, any construction of grue contains, in its first disjunct, a construction that agrees with green on that domain, and in its second, one that agrees with blue on the rest. Its shortest construction is accordingly longer than green's by at least the constructions of blue, of being examined, and of the temporal relation to $t$.

\begin{corollary}\label{c:green}
Green is more natural than grue, and \enquote{green} is more projectible over emeralds than \enquote{grue}.
\end{corollary}

\begin{proof}
The first claim follows from Definition~\ref{d:nat} and the comparison of constructions; the second from the first and \ref{p:track}, since the evidence about examined emeralds is the same for both predicates.
\end{proof}

The argument does not rely on the fact that \enquote{grue} is defined, in English, in terms of \enquote{green}. It relies on Proposition~\ref{p:dup}, which is a claim about the properties and the objects that have them, and which would be true in any language.

\subsection{The symmetry objection}

Goodman's reply to Carnap was that whether a predicate is qualitative depends on the language: in the grue--bleen language, \enquote{green} is defined as \enquote{grue if examined before $t$, bleen otherwise}, and appears positional. EPI-04-b examines this objection in detail. The answer in brief is that Proposition~\ref{p:dup} breaks the symmetry. Whatever language is used, the property denoted by \enquote{green} does not distinguish the duplicates $a$ and $b$, and the property denoted by \enquote{grue} does. The grue--bleen speaker can define green in her vocabulary, but the definition she gives is long in the measure that matters, since its length is computed from the universals, not from her primitives. Definition~\ref{d:nat} was designed to make naturalness language-independent, and it is the language-independence that removes the symmetry Goodman exploited against Carnap. Carnap's qualitativeness requirement, so reconstructed, was on the right track: it failed because it was stated in terms of predicates rather than properties.

\section{Cases}

\subsection{Time-dependent properties that project}

Some properties change at a time and are nonetheless projectible. A sample of a radioactive isotope will have lost half its nuclei after a fixed interval; a caterpillar will be a butterfly within weeks; the leaves of a maple will turn red in autumn. Predicates that encode such changes, such as \enquote{has a half-life of 5730 years}, are projectible, and the reason confirms the account. The temporal element enters their construction through a genuine relation, the elapsed duration since some event involving the object itself, and not through a particular date. Two intrinsic duplicates formed at the same time decay alike, and the property does not distinguish duplicates in the way grue does. A property that changes at a time fixed independently of the object, such as grue, has no such construction, and it is that feature, not temporal change as such, that makes it unnatural.

\citet{Jackson1975} argued that grue is projectible in special circumstances, for example for emeralds that would change colour if examined after $t$. We agree, and the account explains why. If there were a law connecting examination after $t$ with a change of colour, the predicate \enquote{grue}, restricted to such emeralds, would be equivalent to a construction from that law's terms, and the generalization that they are grue would be lawlike by \ref{p:law}. What the riddle needs is not that grue is never projectible but that, absent such a law, it is less projectible than green on the same evidence, which \ref{p:track} delivers.

\subsection{Disjunctive kinds}

Jade, as \citet{Putnam1975} observed, is two minerals, jadeite and nephrite, with different chemical compositions. Generalizations over jade that hold because of chemical composition, such as those concerning hardness or crystal structure, are less projectible from a sample of jade than the corresponding generalizations over jadeite, since the sample may be all jadeite and the unobserved jade nephrite. The kind \emph{jade} is less natural than the kind \emph{jadeite}, being their disjunction, and the projectibility of predicates over it is correspondingly weaker. This is a case in which naturalness of the kind, rather than of the predicate, governs projectibility, and it shows that \ref{p:track}, which holds the kind fixed, captures only part of the dependence. The general claim is that projectibility depends on the naturalness of both the kind and the predicate, as MET-27 suggests in holding that natural kinds are real divisions in nature.

\subsection{Observation-relative predicates}

\enquote{Examined}, \enquote{observed by us} and \enquote{described in the literature} are not projectible from examined to unexamined cases, for the obvious reason that the evidence is selected for them. The account classifies them correctly: their constructions involve the particular observers as parameters, and the resemblance between examined emeralds in respect of having been examined is not a resemblance in any respect that the emeralds' natures determine. Grue inherits its unprojectibility in part from containing such a component.

\section{Objections}

\begin{objection}[Circularity]
We identify the natural properties by finding which predicates support successful inductions. The account explains projectibility by naturalness, but our only access to naturalness is through projectibility, so the explanation is empty \citep{Quine1969}.
\end{objection}

\begin{reply}
The objection conflates the order of explanation with the order of discovery. That we learn which properties are natural partly by learning which generalizations hold does not make naturalness consist in projectibility, any more than the fact that we learn which objects are massive by observing their behaviour makes mass consist in behaviour. The epistemic circle, in which we use induction over natural predicates to discover which predicates are natural, is rule-circular in the sense of EPI-03: the inference uses the projectibility of natural predicates without assuming as a premise which predicates are natural. Such circularity is benign when the rule is reliable, and the present account says when it is: in a world whose laws are stated in natural terms.
\end{reply}

\begin{objection}[Entrenchment suffices]
Goodman's entrenchment theory explains the data without metaphysics. \enquote{Green} has been projected for millennia and \enquote{grue} never; that is all that distinguishes them, and all that is needed to explain our practice \citep{Goodman1955}.
\end{objection}

\begin{reply}
Entrenchment explains our practice; it does not explain why the practice succeeds, and the riddle asks which predicates are projectible, not which ones we project. A community that had projected \enquote{grue} since antiquity would have entrenched it, and after $t$ its predictions would fail. Entrenchment is a good guide to projectibility only if the entrenched predicates are the ones whose projections have succeeded, and they have succeeded, on the present account, because they denote natural properties. Entrenchment is the epistemic trace of naturalness, not a substitute for it.
\end{reply}

\begin{objection}[Priors do the work]
In a Bayesian framework the green hypothesis is confirmed and the grue hypothesis is not because the prior probability of the first is higher. The riddle is solved by priors, and naturalness is idle \citep{Sober1994}.
\end{objection}

\begin{reply}
The Bayesian is right that the difference shows up in the priors, and wrong that it ends there. A prior that favoured grue-like hypotheses would make counter-natural projection rational, and the Bayesian framework by itself gives no reason to prefer one prior to the other. EPI-17 argues that for any body of evidence there is a unique rational credence, and that the rational prior is constrained by the projection of natural rather than gerrymandered predicates. The present account supplies the constraint and its justification: the prior should favour hypotheses in natural terms because, by \ref{p:law}, those are the hypotheses that can be laws. The claim that inductive inference requires substantive background assumptions \citep{Okasha2007} is correct, and naturalness is the assumption.
\end{reply}

\begin{objection}[Why care about naturalness?]
Even if natural properties exist, it is unclear why inquiry should aim at them rather than at gerrymandered ones; naturalness is a metaphysical feature, and nothing about it makes it normatively significant \citep{Dasgupta2018}.
\end{objection}

\begin{reply}
Projectibility supplies the reason. An inquirer who aims at prediction must aim at properties whose distribution in the observed cases is evidence of their distribution in the unobserved ones, and by \ref{p:track} these are the natural properties. The normative significance of naturalness is inherited from that of successful prediction, and it does not depend on any further evaluative fact about the natural properties themselves.
\end{reply}

\section{What the result does not establish}

The result does not establish which properties are natural; that is a question for physics and for the theory of properties of MET-02. It does not establish that our world is one in which natural predicates project successfully, which depends on the world's being governed by laws stated in natural terms; EPI-03 shows that induction is justified in such a world and argues that ours appears to be one. It does not provide a numerical measure of projectibility, since the complexity measure of Definition~\ref{d:nat} is comparative and, for most ordinary properties, unknown in detail. And it does not settle the symmetry objection in its full generality, for which we rely on EPI-04-b.

\section{Conclusion}

\enquote{Green} is projectible and \enquote{grue} is not because green is a natural property and grue is not. Naturalness is measured by the complexity of a property's shortest construction from the sparse universals. Inductive inference projects resemblance, resemblance is the sharing of natural properties, and laws are stated in natural terms; so, given a kind and the evidence, projectibility tracks naturalness. Grue distinguishes intrinsic duplicates and can be constructed only with a particular time as a parameter, by a construction that contains green's; it is therefore less natural, and less projectible, in any language.

\begin{thebibliography}{99}\small
\bibitem[Carnap(1947)]{Carnap1947} Carnap, R. (1947). On the application of inductive logic. \emph{Philosophy and Phenomenological Research}, 8(1), 133--148.
\bibitem[Dasgupta(2018)]{Dasgupta2018} Dasgupta, S. (2018). Realism and the absence of value. \emph{The Philosophical Review}, 127(3), 279--322.
\bibitem[Godfrey-Smith(2003)]{GodfreySmith2003} Godfrey-Smith, P. (2003). Goodman's problem and scientific methodology. \emph{The Journal of Philosophy}, 100(11), 573--590.
\bibitem[Goodman(1946)]{Goodman1946} Goodman, N. (1946). A query on confirmation. \emph{The Journal of Philosophy}, 43(14), 383--385.
\bibitem[Goodman(1955)]{Goodman1955} Goodman, N. (1955). \emph{Fact, Fiction, and Forecast}. Harvard University Press.
\bibitem[Hempel(1945)]{Hempel1945} Hempel, C.~G. (1945). Studies in the logic of confirmation. \emph{Mind}, 54(213), 1--26.
\bibitem[Jackson(1975)]{Jackson1975} Jackson, F. (1975). Grue. \emph{The Journal of Philosophy}, 72(5), 113--131.
\bibitem[Lewis(1983)]{Lewis1983} Lewis, D. (1983). New work for a theory of universals. \emph{Australasian Journal of Philosophy}, 61(4), 343--377.
\bibitem[Lewis(1986)]{Lewis1986} Lewis, D. (1986). \emph{On the Plurality of Worlds}. Blackwell.
\bibitem[Norton(2006)]{Norton2006} Norton, J.~D. (2006). How the formal equivalence of grue and green defeats what is new in the new riddle of induction. \emph{Synthese}, 150(2), 185--207.
\bibitem[Okasha(2007)]{Okasha2007} Okasha, S. (2007). What does Goodman's \enquote{grue} problem really show? \emph{Philosophical Papers}, 36(3), 483--502.
\bibitem[Putnam(1975)]{Putnam1975} Putnam, H. (1975). The meaning of \enquote{meaning}. \emph{Minnesota Studies in the Philosophy of Science}, 7, 131--193.
\bibitem[Quine(1969)]{Quine1969} Quine, W.~V. (1969). Natural kinds. In \emph{Ontological Relativity and Other Essays} (pp.~114--138). Columbia University Press.
\bibitem[Sober(1994)]{Sober1994} Sober, E. (1994). No model, no inference: A Bayesian primer on the grue problem. In D.~Stalker (Ed.), \emph{Grue! The New Riddle of Induction}. Open Court.
\end{thebibliography}

\end{document}
